% titel: Blatt 8, A6: Dipolmomente erklären
% Teil II: der durchdachte Weg.

\subsection*{Gegeben und gesucht}
a) $\mu(\mathrm{OF_2}) = 1.0\cdot10^{-30}\ C\,m$, aber $\mu(\mathrm{BeF_2}) = 0$.
b) $\mu(\mathrm{PF_3}) = 3.4\cdot10^{-30}\ C\,m$, aber $\mu(\mathrm{BF_3}) = 0$.
Erkläre den Unterschied.

\subsection*{a) OF$_2$ vs.\ BeF$_2$}
BeF$_2$: Beryllium bildet mit zwei F eine \textbf{lineare} Struktur. Die
beiden identischen Be--F-Bindungsmomente zeigen in entgegengesetzte
Richtungen
$\Rightarrow$ sie heben sich auf, $\mu = 0$ (trotz polarer Bindung).

OF$_2$: Sauerstoff hat zwei Bindungen und zwei freie Elektronenpaare
$\Rightarrow$ \textbf{gewinkelte} Geometrie (V-förmig, wie im Wasser).
Die beiden O--F-Bindungsmomente addieren sich vektoriell nicht zu Null,
also bleibt ein Restmoment $\mu \approx 1.0\cdot10^{-30}\ C\,m\neq 0$.

\subsection*{b) PF$_3$ vs.\ BF$_3$}
BF$_3$: Bor mit drei F bildet eine \textbf{trigonal-planare} Struktur
($120°$-Winkel). Die drei gleich grossen B--F-Momente in der Ebene heben sich
symmetrisch auf $\Rightarrow \mu = 0$.

PF$_3$: Phosphor hat drei Bindungen und ein freies Elektronenpaar $\Rightarrow$
\textbf{trigonal-pyramidale} Geometrie. Die drei P--F-Momente zeigen nach
einer Seite (das freie Paar "stösst" sie zusammen) und addieren sich zu einem
deutlichen Restmoment $\mu = 3.4\cdot10^{-30}\ C\,m\neq 0$.

\begin{ergebnis}
Die Symmetrie bestimmt das Gesamtdipolmoment: linear (BeF$_2$) und
trigonal-planar (BF$_3$) $\Rightarrow \mu=0$; gewinkelt (OF$_2$) und pyramidal
(PF$_3$) $\Rightarrow \mu\neq 0$.
\end{ergebnis}

\begin{kontrolle}
Bei gleichen Liganden kompensieren sich in den symmetrischen Strukturen alle
Teilmomente. Die nicht-symmetrischen Strukturen von OF$_2$ und PF$_3$
erzeugen ein Restmoment --- genau so erwartet man es nach VSEPR.
\end{kontrolle}

\begin{merksatz}
Diplmoment kommt nicht von einzelnen polaren Bindungen allein, sondern vom
Vektor-Summe aller Beiträge. Symmetrische Moleküle mit gleichen Liganden sind
unpolar, unsymmetrische (gewinkelt, pyramidal, trigonal-pyramidal) sind polar.
\end{merksatz}
